\begin{proof}[Proof of \cref{obj:thm:welfare-identity}]
1.  From the well-formedness of the observed law, extract the measurability of
\(\tau_P\) and the identity \(\tau_P(x)=\mu_1(x)-\mu_0(x)\).  By
\cref{obj:ass:bounded-outcome}, \(|\mu_0(x)|\le1\) and
\(|\mu_1(x)|\le1\) for every \(x\).  Hence, for every \(x\),
\[
|\tau_P(x)|
=
|\mu_1(x)-\mu_0(x)|
\le |\mu_1(x)|+|\mu_0(x)|
\le 2 .
\]
% lean: regret_eq_disagreement_integral

2.  The set \(\{x:0\le\tau_P(x)\}\) is measurable because \(\tau_P\) is
measurable, and the set \(\{x:\pi(x)=1\}\) is measurable because \(\pi\) is
measurable.  Therefore the two integrands
\[
x\mapsto
\begin{cases}
\tau_P(x),&0\le \tau_P(x),\\
0,&\text{otherwise},
\end{cases}
\qquad
x\mapsto
\begin{cases}
\tau_P(x),&\pi(x)=1,\\
0,&\text{otherwise},
\end{cases}
\]
are measurable and, by the preceding bound \(|\tau_P|\le2\), integrable with
respect to the probability measure \(P_X\).  This is the integrability bookkeeping needed to
combine the two welfare integrals into one integral of a difference.
% lean: regret_eq_disagreement_integral

3.  By \cref{obj:def:welfare-regret}, \(V_P(\pi)=\int\mathbf 1\{\pi(x)=1\}\tau_P(x)\,dP_X(x)\) and \(\pi^\star_P(x)=\mathbf 1\{\tau_P(x)\ge0\}\), so \(\mathbf 1\{\pi^\star_P(x)=1\}=\mathbf 1\{0\le\tau_P(x)\}\).  Expanding \(R_P(\pi)=V_P(\pi^\star_P)-V_P(\pi)\) therefore gives
\[
R_P(\pi)
=
\int
\left[
\mathbf 1\{0\le\tau_P(x)\}\tau_P(x)
-
\mathbf 1\{\pi(x)=1\}\tau_P(x)
\right]\,dP_X(x).
\]
The integrability of Step 2 justifies writing the difference of the two integrals as the integral of the pointwise
difference.
% lean: regret_eq_disagreement_integral

4.  It remains only to identify the pointwise difference.  Fix \(x\).  If
\(\tau_P(x)\ge0\), then \(\pi^\star_P(x)=1\).  When \(\pi(x)=1\), both terms equal \(\tau_P(x)\) and the difference is
zero, and indeed there is no disagreement; when \(\pi(x)\ne1\), the difference is \(\tau_P(x)=|\tau_P(x)|\), and there is
disagreement.  If \(\tau_P(x)<0\), then \(\pi^\star_P(x)=0\), so the first term is zero.
When \(\pi(x)=1\), the difference is \(-\tau_P(x)=|\tau_P(x)|\), and there is disagreement; when
\(\pi(x)\ne1\), both terms are zero and there is no disagreement.  Thus, for every \(x\),
\[
\mathbf 1\{0\le\tau_P(x)\}\tau_P(x)
-
\mathbf 1\{\pi(x)=1\}\tau_P(x)
=
|\tau_P(x)|\,\mathbf 1\{\pi(x)\ne\pi^\star_P(x)\}.
\]
% lean: regret_eq_disagreement_integral

5.  Substituting this pointwise identity into the integral yields
\[
R_P(\pi)
=
\int_{\mathcal X}
|\tau_P(x)|\,\mathbf 1\{\pi(x)\ne\pi^\star_P(x)\}\,dP_X(x)
=
E_P\!\left[|\tau_P(X)|\,\mathbf 1\{\pi(X)\ne\pi^\star_P(X)\}\right],
\]
the second equality because \(P_X\) is the covariate marginal of \(P\) and the integrand is a function of \(X\) alone.  This is the disagreement-integral representation of the regret.
% lean: regret_eq_disagreement_integral
\end{proof}
